% Generated draft: requires mathematical and source review.
\documentclass[11pt]{article}
\usepackage[margin=25mm]{geometry}
\usepackage{fontspec}
\setmainfont{Latin Modern Roman}
\usepackage{amsmath,amssymb,mathtools,unicode-math}
\setmathfont{Latin Modern Math}
\usepackage{enumitem,needspace,xurl,hyperref}
\setlength{\parindent}{0pt}
\setlength{\parskip}{5pt}
\setlength{\emergencystretch}{3em}
\newcommand{\N}{\mathbb{N}}
\newcommand{\Z}{\mathbb{Z}}
\newcommand{\Q}{\mathbb{Q}}
\newcommand{\R}{\mathbb{R}}
\newcommand{\C}{\mathbb{C}}
\newcommand{\F}{\mathbb{F}}
\newcommand{\A}{\mathbb{A}}
\newcommand{\PP}{\mathbb{P}}
\newcommand{\E}{\mathbb{E}}
\newcommand{\eps}{\varepsilon}
\newcommand{\dd}{\,\mathrm{d}}
\begin{document}
\section*{UnsolvedMath ID: 20003292}
\textbf{Source problem code:} catalogue entry
\subsection{Definitions and mathematical statement}
This OCR fragment is not a standalone Conjecture 20. The supplied assessment identifies it as truncated status prose following Conjecture 21; it ends mid-sentence at ``Furthermore, the statement of''. It reports a proof in [136] under $c_1M^{-\nu}\leq y\leq c_2M^{-\nu}$ for subcritical $\nu<2$, but does not reproduce the full theorem or continuation.

Neighboring local records place this in a lattice surface $X=\Gamma\backslash\mathbb H$ with a period-one cusp, a closed horocycle $u_y(t)$ of length $y^{-1}$, and sampling at a badly approximable rotation $t=\alpha m$. The separately owned Conjecture 21 of J. Marklof and A. Str\"ombergsson asks for uniform convergence of the sampled average of a test function to its Liouville mean throughout the critical band $c_1M^{-2}\leq y\leq c_2M^{-2}$, with a specified but unreproduced cusp-growth condition. The present record is only status context, not that full conjecture.

The aid's abstract audit fixes nonnegative errors $E_M(y)$, $0<c_1<c_2$, and $M\to\infty$. A genuine one-sided theorem asserting $E_M(y)\to0$ whenever simultaneously $y\to0$ and $y\geq cM^{-\nu_0}$ transfers to every fixed band with $0<\nu<\nu_0$, because $c_1M^{-\nu}\geq cM^{-\nu_0}$ eventually. This uses a uniform/two-parameter one-sided hypothesis. Convergence confined to one fixed comparable band does not supply that hypothesis or imply convergence at smaller $y$.

In particular, even uniform convergence on every fixed subcritical band does not prove the exponent-$2$ endpoint. The explicit countermodel is
\[
 E_M(y)=e^{-M^2y}.
\]
For any fixed $0<\nu<2$,
\[
 \sup_{c_1M^{-\nu}\leq y\leq c_2M^{-\nu}}E_M(y)
       =e^{-c_1M^{2-\nu}}\longrightarrow0,
\]
whereas on the endpoint band the supremum is $e^{-c_1}$ and $E_M(cM^{-2})=e^{-c}$. This refutes an inference between quantified estimates, not the actual horocycle conjecture.

For a quantitative sufficient interpolation, suppose on the required region
\[
 E_M(y)\leq r_M+C(M^2y)^{-\delta},\qquad
 r_M\to0,\quad C,\delta>0,
\]
with constants independent of $M,y$. For a moving exponent $\nu_M<2$ tending to $2$ and $y\geq c_1M^{-\nu_M}$, this gives
\[
 E_M(y)\leq r_M+C c_1^{-\delta}
   \exp\bigl(-\delta(2-\nu_M)\log M\bigr).
\]
Thus $(2-\nu_M)\log M\to\infty$ is a sufficient condition for uniform decay in the corresponding bands. Failure of that condition means this bound alone gives no decay, not that the true dynamical error must fail to converge.

The aid separately reports a strict cusp-growth boundary $\gamma<3/4$ at $\nu=2$. Its function class, height normalization and underlying inequality are not supplied here, so this is retained as a reported threshold needing the full [136] setup; no theorem at equality $\gamma=3/4$ or for every continuous cusp-growing function is asserted. Cusp integrability and passage to the critical sampling scale are separate requirements.
\subsection{Short English statement}
This is truncated status prose, not a new conjecture. The supplied audit explains why convergence for each fixed exponent below $2$ does not reach the critical exponent: $e^{-M^2y}$ is an exact countermodel to that inference. A quantitative error estimate can approach the endpoint only with an explicit moving-exponent condition, and cusp-growth restrictions require separate control.
\subsection{Sources}
\begin{itemize}
\item Local UnsolvedMath registry, permanent ID $20003292$, source code catalogue entry; canonical location \texttt{aim-other-notes.json}, \texttt{notes[98]}. Neighboring records $20003290,20003291$ supply the horocycle context and Conjecture 21 attribution to J. Marklof and A. Str\"ombergsson.
\item AIM \emph{Emerging applications of measure rigidity}, truncated status text following Conjecture 21, \url{https://aimath.org/WWN/measrigid/measrigid.pdf}. Source label [136] is supplied without a full bibliography entry; the continuation is missing.
\item Research attempt 1 is marked context-only/not-a-problem, proved, and low novelty confidence, despite the local conjecture/open labels. It explicitly warns that the aid is machine generated, not peer reviewed, and requires independent verification before citation.
\end{itemize}
\end{document}
